\documentclass[11pt,a4paper]{article}

\usepackage[utf8]{inputenc}
\usepackage[T1]{fontenc}
\usepackage{lmodern}
\usepackage[spanish,es-noshorthands]{babel}
\usepackage[margin=2.5cm,top=2.6cm,bottom=2.6cm]{geometry}
\usepackage{amsmath,amssymb,amsthm}
\usepackage{booktabs,array}
\usepackage[table]{xcolor}
\usepackage{tikz}
\usetikzlibrary{calc,positioning,arrows.meta}
\usepackage{fancyhdr}
\usepackage{titlesec}
\usepackage{tocloft}
\usepackage[font=small,labelfont=bf]{caption}
\usepackage{enumitem}
\usepackage{algorithm}
\usepackage{algpseudocode}
\usepackage[colorlinks=true,linkcolor=acento,citecolor=acento,urlcolor=acento]{hyperref}

% ---------- colores ----------
\definecolor{fondo}{HTML}{263A42}
\definecolor{franja}{HTML}{7D969E}
\definecolor{acento}{HTML}{2F5F73}
\definecolor{celda}{HTML}{CFDAE0}
\definecolor{gris}{HTML}{6B7C84}

\AtBeginDocument{\renewcommand{\tablename}{Tabla}\renewcommand{\contentsname}{Contenido}\renewcommand{\figurename}{Figura}\renewcommand{\proofname}{Demostración}\renewcommand{\listalgorithmname}{Algoritmos}}
\floatname{algorithm}{Algoritmo}
\algrenewcommand\algorithmicfor{\textbf{para}}
\algrenewcommand\algorithmicdo{\textbf{hacer}}
\algrenewcommand\algorithmicif{\textbf{si}}
\algrenewcommand\algorithmicthen{\textbf{entonces}}
\algrenewcommand\algorithmicend{\textbf{fin}}
\algrenewcommand\algorithmicreturn{\textbf{devolver}}

% ---------- macros ----------
\newcommand{\MochilaDec}{\textsc{Mochila-Dec}}
\newcommand{\AcelDec}{\textsc{Acel-Dec}}
\newcommand{\DP}[2]{\mathrm{DP}[#1][#2]}
\newcommand{\sh}{\cellcolor{celda}}

% ---------- teoremas ----------
\newtheoremstyle{plano}{8pt}{8pt}{\upshape}{}{\bfseries}{.}{ }{\thmname{#1}\thmnumber{ #2}\thmnote{ (#3)}}
\theoremstyle{plano}
\newtheorem{definicion}{Definición}[section]
\newtheorem{proposicion}{Proposición}[section]
\renewcommand{\qedsymbol}{$\blacksquare$}

% ---------- encabezados ----------
\pagestyle{fancy}
\fancyhf{}
\lhead{Análisis y Diseño de Algoritmos \,$\cdot$\, 2026-B}
\rhead{Grupo 3 \,$\cdot$\, Primera entrega}
\cfoot{\thepage}
\renewcommand{\headrulewidth}{0.4pt}

\titleformat{\section}{\Large\bfseries}{\thesection}{1em}{}
\titleformat{\subsection}{\large\bfseries}{\thesubsection}{1em}{}
\titlespacing*{\section}{0pt}{22pt}{10pt}
\titlespacing*{\subsection}{0pt}{16pt}{7pt}
\setlength{\parindent}{0pt}
\setlength{\parskip}{6pt}

% ---------- índice ----------
\renewcommand{\cftsecfont}{\bfseries\color{acento}}
\renewcommand{\cftsecpagefont}{\bfseries}
\renewcommand{\cftsubsecfont}{\color{acento}}
\renewcommand{\cftsubsecpagefont}{\color{acento}}
\setlength{\cftbeforesecskip}{9pt}
\setlength{\cftbeforesubsecskip}{3pt}
\renewcommand{\cftdotsep}{2}

\setlist[itemize]{itemsep=2pt,topsep=2pt}

\begin{document}

% ======================= PORTADA =======================
\begin{titlepage}
\thispagestyle{empty}
\begin{tikzpicture}[remember picture,overlay]
  \fill[fondo] (current page.north west) rectangle ($(current page.north east)+(0,-9.6cm)$);
  \fill[franja] ($(current page.north west)+(0,-9.6cm)$) rectangle ($(current page.north east)+(0,-9.9cm)$);
  \node[anchor=north west,text=white,inner sep=0] at ($(current page.north west)+(2.5cm,-1.9cm)$) {%
    \begin{minipage}{14cm}
      {\scshape Universidad de Caldas}\\[2pt]
      {\small Análisis y Diseño de Algoritmos \,$\cdot$\, 2026-B}\\[26pt]
      {\fontsize{26}{33}\selectfont\bfseries Aceleración de la cadena\\bajo presupuesto}\\[8pt]
      {\normalsize Primera entrega \,$\cdot$\, Subproyecto 3 (Grupo 3)}
    \end{minipage}};
\end{tikzpicture}
\vspace*{10.2cm}

\begin{tabular}{@{}p{3.6cm}l@{}}
  \textbf{\color{acento}Integrantes} & Rafael Medina\\[6pt]
   & Isabela Guerrero\\[6pt]
  \textbf{\color{acento}Asignatura} & Análisis y Diseño de Algoritmos, 2026-B\\[6pt]
  \textbf{\color{acento}Institución} & Universidad de Caldas\\[6pt]
  \textbf{\color{acento}Entrega} & Viernes 6 de noviembre de 2026\\[6pt]
  \textbf{\color{acento}Repositorio} & {\small\ttfamily\href{https://github.com/Mandara2/ProyectoADAG3.git}{https://github.com/Mandara2/ProyectoADAG3.git}}
\end{tabular}
\end{titlepage}

\setcounter{page}{1}
\tableofcontents
\newpage

% ======================= 1 ALCANCE =======================
\section{Alcance}
\label{sec:alcance}

Este informe resuelve la primera parte del Subproyecto 3: decidir qué etapas de una cadena conviene acelerar con un presupuesto limitado. La cobertura mínima de auditoría (la segunda parte) queda fuera de esta entrega, por indicación de la profesora.

El código está en el repositorio \url{https://github.com/Mandara2/ProyectoADAG3.git}.

% ======================= 2 PROBLEMA =======================
\section{El problema y su formalización}

\textbf{En palabras simples.} Una solicitud pasa por una fila de etapas, una tras otra. Cada etapa tarda un tiempo y se puede acelerar pagando un costo, lo que le quita algo de tiempo. Hay un presupuesto total y no se puede gastar más de eso. La pregunta es qué etapas conviene acelerar para que la cadena quede lo más corta posible.

Como las etapas van en fila, todas cuentan para el tiempo total. Por eso basta con sumar las duraciones, y el problema se reduce a esto: \textbf{ganar la mayor cantidad de tiempo posible sin pasarse del presupuesto}.

\subsection{Tamaño de la entrada}

Se mide con dos números:
\begin{itemize}
  \item $n$, el número de etapas. Fija cuántas decisiones hay que tomar (acelerar o no cada etapa).
  \item $B$, el presupuesto entero. Fija cuántos niveles de gasto hay que considerar.
\end{itemize}
La longitud real de la entrada en bits es
\begin{equation}
  L = O\Big(\sum_{i=1}^{n}\big(\log(t_i+1) + \log(c_i+1) + \log(d_i+1)\big) + \log(B+1)\Big).
\end{equation}
Se usa $\log(x+1)$ porque un número $x\ge 0$ ocupa $\lceil\log_2(x+1)\rceil$ bits, incluso si vale 0.

\subsection{Entrada, salida, restricción y objetivo}

\begin{definicion}[Aceleración de la cadena]
\textbf{Entrada.} Un entero $n\ge 0$, un entero $B\ge 0$ y, para cada etapa $i\in\{1,\dots,n\}$, tres enteros no negativos:
\begin{itemize}
  \item $t_i$: duración normal de la etapa;
  \item $c_i$: costo de acelerarla ($c_i\ge 0$; un costo negativo haría que la tabla de la sección~\ref{sec:sol} se saliera de sus columnas);
  \item $d_i$: tiempo que se ahorra si se acelera, con $0\le d_i\le t_i$ (no se puede ahorrar más de lo que dura).
\end{itemize}
\textbf{Salida.} Un conjunto $A\subseteq\{1,\dots,n\}$ con las etapas que se aceleran.

\textbf{Restricción.} El costo total no puede superar el presupuesto:
\begin{equation}
  \sum_{i\in A} c_i \le B.
\end{equation}
\textbf{Objetivo.} Minimizar la duración total, que es la suma de todas las duraciones menos lo que se ahorra en las etapas aceleradas:
\begin{equation}
  \min_{A}\; D(A)=\sum_{i=1}^{n} t_i-\sum_{i\in A} d_i.
\end{equation}
\end{definicion}
La primera suma no depende de lo que se elija, así que minimizar la duración es lo mismo que maximizar el tiempo ahorrado:
\begin{equation}
  \max_{A}\;\sum_{i\in A} d_i \qquad\text{sujeto a}\qquad \sum_{i\in A} c_i\le B.
\end{equation}
En resumen, es el clásico problema de la mochila: el presupuesto es la capacidad, el costo es el peso y el tiempo ahorrado es el valor.

% ======================= 3 MARCO TEORICO =======================
\section{Marco teórico: P, NP y NP-dificultad}

\textbf{Idea general.} Hay problemas que se resuelven rápido, y otros en los que comprobar una respuesta es fácil pero encontrarla parece muy costoso. Un problema es NP-difícil cuando es, como mínimo, tan difícil como cualquiera de esos problemas «duros». Esta sección define los conceptos que se usan en la prueba de la sección~\ref{sec:np}.

\begin{table}[h]
\centering\small
\renewcommand{\arraystretch}{1.15}
\begin{tabular}{@{}>{\bfseries\raggedright}p{2.5cm}p{4.3cm}p{6.3cm}@{}}
\toprule
\normalfont\bfseries Concepto & En palabras simples & Definición formal \\
\midrule
Problema de decisión & Una pregunta con respuesta sí o no. & Las clases de complejidad se definen sobre ellos. Un problema de optimización, como el de este proyecto, se asocia con una versión de decisión que agrega un umbral: ¿existe una solución con duración de a lo sumo $T$? \\[4pt]
Tiempo polinomial & «Rápido»: al duplicar la entrada, el tiempo se multiplica por una cantidad fija (por 4 si es $O(L^2)$). En uno exponencial, en cambio, el tiempo se eleva al cuadrado. & Un algoritmo es polinomial si existe una constante $k$ tal que su número de pasos es $O(L^k)$, con $L$ la longitud en bits de la entrada. \\[4pt]
Clase P & Problemas que se resuelven rápido. & Problemas de decisión que se pueden resolver en tiempo polinomial. \\[4pt]
Clase NP & Problemas donde, si alguien te da la respuesta, puedes comprobarla rápido (como revisar un sudoku ya resuelto). & Problemas de decisión en los que, cuando la respuesta es sí, existe un certificado de tamaño polinomial que se verifica en tiempo polinomial. Todo problema de P está en NP. \\[4pt]
NP-difícil & Al menos tan difícil como cualquier problema de NP. & $X$ es NP-difícil si todo problema de NP se reduce a $X$ en tiempo polinomial. $X$ no tiene que estar en NP. \\[4pt]
NP-completo & Difícil y, además, fácil de comprobar. & NP-difícil y además en NP. \\
\bottomrule
\end{tabular}
\caption{Conceptos de complejidad usados en el informe.}
\end{table}

\begin{definicion}[Reducción polinomial]
Un problema $A$ se reduce a un problema $B$, y se escribe $A\le_p B$, si existe una función $f$ computable en tiempo polinomial tal que
\begin{equation}
  x \text{ es un caso sí de } A \iff f(x) \text{ es un caso sí de } B.
\end{equation}
\end{definicion}
La idea es esta: si supiéramos resolver $B$ rápido, resolveríamos $A$ rápido, porque bastaría con convertir el caso y resolver el resultado. Dicho al revés, $B$ es al menos tan difícil como $A$. Además, las reducciones se pueden encadenar: si $A\le_p B$ y $B\le_p C$, entonces $A\le_p C$.

\textbf{Cómo se prueba en la práctica.} Cook (1971) mostró que SAT (el problema de satisfacibilidad booleana: dada una fórmula lógica, decidir si existe una asignación de verdadero/falso a sus variables que la haga verdadera) es NP-completo, y Karp (1972) probó lo mismo para otros problemas, entre ellos Mochila 0/1. Para probar que un problema nuevo $X$ es NP-difícil no hace falta reducir todo NP a $X$: basta con reducir a $X$ un problema que ya se sabe NP-difícil. Por transitividad, todo NP queda reducido a $X$.

\textbf{Por qué importa.} No se sabe si P $=$ NP, pero se conjetura que no. Si P $\ne$ NP, ningún problema NP-difícil tiene un algoritmo polinomial. Por eso, probar que este problema es NP-difícil justifica usar un algoritmo exacto pseudo-polinomial (sección~\ref{sec:sol}) en lugar de buscar uno polinomial.

\textbf{NP-difícil en sentido débil.} Es un problema NP-difícil que tiene un algoritmo pseudo-polinomial: su tiempo es polinomial en el valor numérico de los datos, pero no en su longitud en bits. Mochila 0/1 es de este tipo. Nuestro problema también lo es, por dos razones distintas: es NP-difícil por la reducción de la sección~\ref{sec:np}, y tiene un algoritmo pseudo-polinomial propio, la programación dinámica de la sección~\ref{sec:sol}.

% ======================= 4 NP-DUREZA =======================
\section{Prueba de NP-dureza}
\label{sec:np}

\textit{Idea.} Acelerar la cadena es el problema de la mochila con otro disfraz. Cada objeto de la mochila se convierte en una etapa: el peso pasa a ser el costo de acelerarla, y el valor pasa a ser el tiempo que se ahorra. Si supiéramos resolver rápido la aceleración, resolveríamos rápido la mochila, y como la mochila es NP-difícil, la aceleración también lo es. La reducción va de Mochila hacia la aceleración, que es la dirección que exige la Proposición 4.2.1 del enunciado.

\subsection{Las dos versiones de decisión}

\begin{definicion}[\textsc{Acel-Dec}, el problema nuevo]
Dadas las etapas, el presupuesto $B$ y un entero $T$: ¿existe $A$ con costo $\le B$ y duración $\le T$?
\begin{equation}
  \exists A\subseteq\{1,\dots,n\}:\; \sum_{i\in A} c_i\le B \;\wedge\; \sum_{i=1}^{n} t_i-\sum_{i\in A} d_i\le T.
\end{equation}
\end{definicion}
\begin{definicion}[\textsc{Mochila-Dec}, el problema de partida]
Dados objetos con pesos $w_i$ y valores $v_i$ enteros no negativos, una capacidad entera $W\ge 0$ y un umbral entero $V\ge 0$: ¿existe $S$ con peso $\le W$ y valor $\ge V$? Es NP-completo (Karp, 1972).
\begin{equation}
  \exists S\subseteq\{1,\dots,n\}:\; \sum_{i\in S} w_i\le W \;\wedge\; \sum_{i\in S} v_i\ge V.
\end{equation}
\end{definicion}

\subsection{Construcción: el «traductor»}

Para cada objeto $i$ se crea una etapa $i$, y se fijan $B$ y $T$ así:
\begin{equation}
  c_i=w_i,\qquad d_i=v_i,\qquad t_i=v_i,\qquad B=W,\qquad T=\sum_{i=1}^{n} v_i-V.\label{eq:trad}
\end{equation}

\begin{table}[h]
\centering\small
\renewcommand{\arraystretch}{1.1}
\begin{tabular}{@{}ll@{}}
\toprule
\bfseries Mochila & \bfseries Aceleración \\
\midrule
Peso del objeto $w_i$ & Costo de acelerar $c_i$ \\
Valor del objeto $v_i$ & Tiempo ahorrado $d_i$ \\
Capacidad $W$ & Presupuesto $B$ \\
Valor mínimo $V$ & Duración máxima $T=\sum v_i-V$ \\
\bottomrule
\end{tabular}
\caption{Diccionario de la reducción.}
\end{table}

La elección $t_i=v_i$ garantiza $0\le d_i\le t_i$, así que el resultado es un caso válido.

\textbf{Así se ve la equivalencia.} Tomamos el ejemplo de tres etapas con los mismos costos y ahorros del caso \texttt{ejemplo\_tres\_etapas} ($c=4,3,5$ y $d=3,2,4$), con $t=d$ como pide la construcción, $B=10$ y $T=9-7=2$.

% ---------- Figura 1 ----------
\begin{figure}[t]
\centering
\begin{tikzpicture}[x=1cm,y=1cm,font=\footnotesize]
  \tikzset{
    sel/.style={draw=fondo,fill=celda!70,rounded corners=2pt,thick,minimum width=6.4cm,minimum height=0.95cm,inner sep=4pt},
    nosel/.style={draw=gray!50,fill=white,rounded corners=2pt,minimum width=6.4cm,minimum height=0.95cm,inner sep=4pt},
    flecha/.style={-{Stealth[length=2mm]},gray!70}
  }
  \node[anchor=west,font=\small\bfseries] at (0,0.8) {Mochila 0/1};
  \node[anchor=west,font=\small\bfseries] at (9.4,0.8) {Aceleración de la cadena};
  \node[gray] at (7.6,0.8) {\scriptsize $\to$ O($n$)};
  % objetos
  \foreach \i/\y/\w/\v/\s/\txtL/\txtR in {1/0/4/3/sel/{va en la mochila}/{se acelera},2/-1.35/3/2/nosel/{se queda fuera}/{no se acelera},3/-2.7/5/4/sel/{va en la mochila}/{se acelera}}{
    \node[\s,anchor=west] (o\i) at (0,\y) {};
    \node[anchor=north west,font=\footnotesize\bfseries] at ($(o\i.north west)+(0.15,-0.1)$) {Objeto \i};
    \node[anchor=north east,font=\scriptsize] at ($(o\i.north east)+(-0.15,-0.1)$) {\txtL};
    \node[anchor=south west,font=\scriptsize] at ($(o\i.south west)+(0.15,0.1)$) {$w_\i=\w,\ v_\i=\v$};
    \node[\s,anchor=west] (e\i) at (9.4,\y) {};
    \node[anchor=north west,font=\footnotesize\bfseries] at ($(e\i.north west)+(0.15,-0.1)$) {Etapa \i};
    \node[anchor=north east,font=\scriptsize] at ($(e\i.north east)+(-0.15,-0.1)$) {\txtR};
    \node[anchor=south west,font=\scriptsize] at ($(e\i.south west)+(0.15,0.1)$) {$c_\i=\w,\ d_\i=\v,\ t_\i=\v$};
    \draw[flecha] (o\i.east) -- (e\i.west);
  }
  % barras 1: peso / costo
  \node[anchor=west,font=\small\bfseries] at (0,-4.35) {Peso total $\le$ capacidad $W=10$};
  \node[anchor=west,font=\small\bfseries] at (9.4,-4.35) {Costo total $\le$ presupuesto $B=10$};
  \draw[dashed,gray] (0,-4.65) rectangle (5.0,-5.15);
  \draw[fill=celda!70,draw=fondo] (0,-4.65) rectangle (2.0,-5.15) node[midway]{\scriptsize obj. 1};
  \draw[fill=celda!70,draw=fondo] (2.0,-4.65) rectangle (4.5,-5.15) node[midway]{\scriptsize obj. 3};
  \node[anchor=west,gray,font=\scriptsize] at (0,-5.45) {peso $4+5=9\le W=10$};
  \draw[dashed,gray] (9.4,-4.65) rectangle (14.4,-5.15);
  \draw[fill=celda!70,draw=fondo] (9.4,-4.65) rectangle (11.4,-5.15) node[midway]{\scriptsize etapa 1};
  \draw[fill=celda!70,draw=fondo] (11.4,-4.65) rectangle (13.9,-5.15) node[midway]{\scriptsize etapa 3};
  \node[anchor=west,gray,font=\scriptsize] at (9.4,-5.45) {costo $4+5=9\le B=10$};
  \draw[{Stealth[length=2mm]}-{Stealth[length=2mm]},gray] (5.4,-4.9) -- (8.9,-4.9);
  \node[gray,font=\scriptsize,above] at (7.15,-4.9) {equivale a};
  % barras 2: valor / duración
  \node[anchor=west,font=\small\bfseries] at (0,-6.15) {Valor total $\ge$ umbral $V=7$};
  \node[anchor=west,font=\small\bfseries] at (9.4,-6.15) {Duración final $\le T=2$};
  \draw[fill=celda!70,draw=fondo] (0,-6.45) rectangle (1.5,-6.95) node[midway]{\scriptsize obj. 1};
  \draw[draw=gray!50] (1.5,-6.45) rectangle (2.5,-6.95) node[midway,gray]{\scriptsize obj. 2};
  \draw[fill=celda!70,draw=fondo] (2.5,-6.45) rectangle (4.5,-6.95) node[midway]{\scriptsize obj. 3};
  \node[anchor=west,gray,font=\scriptsize] at (0,-7.25) {valor $3+4=7\ge V=7$};
  \draw[fill=celda!70,draw=fondo,dashed,thick] (9.4,-6.45) rectangle (10.9,-6.95) node[midway]{\scriptsize et. 1};
  \draw[draw=gray!50,dashed] (10.9,-6.45) rectangle (11.9,-6.95) node[midway,gray]{\scriptsize et. 2};
  \draw[fill=celda!70,draw=fondo,dashed,thick] (11.9,-6.45) rectangle (13.9,-6.95) node[midway]{\scriptsize et. 3};
  \node[anchor=west,gray,font=\scriptsize] at (9.4,-7.25) {$9-7=2\le T=2$ (punteado = tiempo ahorrado)};
  \draw[{Stealth[length=2mm]}-{Stealth[length=2mm]},gray] (5.4,-6.7) -- (8.9,-6.7);
  \node[gray,font=\scriptsize,above] at (7.15,-6.7) {equivale a};
\end{tikzpicture}
\caption{Equivalencia entre la mochila y la aceleración con el ejemplo de tres objetos/etapas. Los costos son los pesos y los ahorros son los valores, así que la misma elección sirve en ambos problemas.}
\label{fig:equiv}
\end{figure}

La figura~\ref{fig:equiv} se lee así: a la izquierda cada objeto se vuelve una etapa a la derecha, con los mismos números. Las dos barras de arriba muestran que «no pasarse de la capacidad» y «no pasarse del presupuesto» son la misma condición. Las dos de abajo muestran que «llegar al valor $V$» y «dejar la cadena en $T$ o menos» también lo son: lo que la mochila gana en valor, la cadena lo gana en tiempo ahorrado.

\textbf{Tiempo de la construcción.} Se recorren los $n$ objetos una vez y se calcula una suma: $O(n)$ operaciones aritméticas sobre números de la entrada. Es polinomial en la longitud de la entrada.

\subsection{Equivalencia: la respuesta es la misma en ambos problemas}

\begin{proposicion}[NP-dureza]
$\MochilaDec\le_p\AcelDec$. En consecuencia, \AcelDec{} es NP-difícil.
\end{proposicion}
\begin{proof}
Por construcción (\ref{eq:trad}), $\sum t_i=\sum v_i$ y, para cualquier conjunto $X$, $\sum_{i\in X} c_i=\sum_{i\in X} w_i$ y $\sum_{i\in X} d_i=\sum_{i\in X} v_i$.

$(\Rightarrow)$ \textit{Si la mochila tiene solución, la aceleración también.} Sea $S$ una solución de la mochila y se toma $A=S$. El costo cumple la restricción,
\begin{equation}
  \sum_{i\in A} c_i=\sum_{i\in S} w_i\le W=B,
\end{equation}
y la duración no pasa de $T$:
\begin{equation}
  \sum_{i=1}^{n} t_i-\sum_{i\in A} d_i=\sum_{i=1}^{n} v_i-\sum_{i\in S} v_i\le\sum_{i=1}^{n} v_i-V=T.
\end{equation}
$(\Leftarrow)$ \textit{Si la aceleración tiene solución, la mochila también.} Sea $A$ una solución de \AcelDec{} y se toma $S=A$. El peso cumple $\sum w_i=\sum c_i\le B=W$. De la cota de duración se obtiene el valor:
\begin{equation}
  \sum_{i=1}^{n} v_i-\sum_{i\in A} v_i\le\sum_{i=1}^{n} v_i-V \implies \sum_{i\in S} v_i\ge V.
\end{equation}
Como la construcción toma tiempo $O(n)$, la reducción es polinomial.
\end{proof}

\textbf{Conclusión.} \MochilaDec{} se reduce a \AcelDec{} en tiempo polinomial, así que \AcelDec{} es NP-difícil. Además está en NP: el certificado es $A$ y se verifica sumando en $O(n)$. Por tanto, \AcelDec{} es NP-completo.

La versión de optimización también es NP-difícil: un algoritmo polinomial para ella respondería \AcelDec{} comparando la duración óptima con $T$. Como la cadena en serie es un caso particular de un flujo general, el problema general es al menos igual de difícil.

% ======================= 5 SOLUCION =======================
\section{Algoritmos de solución: fuerza bruta y programación dinámica}
\label{sec:sol}

\textit{Idea.} La solución directa es probar todas las combinaciones de etapas, pero son $2^n$, un número inmenso. En cambio, la programación dinámica resuelve primero problemas pequeños y reutiliza sus respuestas. Se arma una tabla que responde: «si solo puedo usar las primeras $i$ etapas y gastar como máximo $b$, ¿cuánto tiempo puedo ahorrar como máximo?». Es el Algoritmo 4.2.1 del enunciado y da la respuesta exacta en $O(n\cdot B)$.

Como la aceleración es la mochila con otros nombres (sección~\ref{sec:np}), los mismos dos algoritmos resuelven ambos problemas. Cada uno se presenta primero para la mochila y después para su equivalente en la aceleración. Para la programación dinámica se entregan dos tablas: una con los \textbf{óptimos} y otra que muestra \textbf{cómo se llegó} a ellos.

\subsection{Fuerza bruta}

\textbf{Mochila.} Se prueban los $2^n$ subconjuntos. Cada uno se codifica con un número \textit{mascara} entre $0$ y $2^n-1$, cuyo bit $i$ indica si el objeto $i$ entra. Para cada uno se suman peso y valor, y se guarda el mejor que cabe.

\begin{algorithm}[h]
\caption{\textsc{Mochila-FB}$(w,v,W,n)$}
\begin{algorithmic}[1]
\State $mejor\gets 0$;\quad $S^*\gets\emptyset$
\For{$mascara\gets 0$ \textbf{hasta} $2^n-1$}
  \State $peso\gets 0$;\quad $valor\gets 0$
  \For{$i\gets 1$ \textbf{hasta} $n$}
    \If{el bit $i$ de $mascara$ vale 1}
      \State $peso\gets peso+w_i$;\quad $valor\gets valor+v_i$
    \EndIf
  \EndFor
  \If{$peso\le W$ \textbf{y} $valor>mejor$}
    \State $mejor\gets valor$;\quad $S^*\gets\{\,i : \text{el bit } i \text{ de } mascara \text{ vale } 1\,\}$
  \EndIf
\EndFor
\State \Return $(mejor,S^*)$
\end{algorithmic}
\end{algorithm}

\textbf{Aceleración (equivalente).} Es el mismo algoritmo con $w_i\to c_i$, $v_i\to d_i$ y $W\to B$, más un último paso que convierte el ahorro máximo en la duración mínima de la cadena.

\begin{algorithm}[h]
\caption{\textsc{Aceleración-FB}$(t,c,d,B,n)$}
\begin{algorithmic}[1]
\State $ahorro^*\gets 0$;\quad $A^*\gets\emptyset$
\For{$mascara\gets 0$ \textbf{hasta} $2^n-1$}
  \State $costo\gets 0$;\quad $ahorro\gets 0$
  \For{$i\gets 1$ \textbf{hasta} $n$}
    \If{el bit $i$ de $mascara$ vale 1}
      \State $costo\gets costo+c_i$;\quad $ahorro\gets ahorro+d_i$
    \EndIf
  \EndFor
  \If{$costo\le B$ \textbf{y} $ahorro>ahorro^*$}
    \State $ahorro^*\gets ahorro$;\quad $A^*\gets\{\,i : \text{el bit } i \text{ de } mascara \text{ vale } 1\,\}$
  \EndIf
\EndFor
\State \Return $\bigl(\sum_i t_i-ahorro^*,\;A^*\bigr)$
\end{algorithmic}
\end{algorithm}

En el ejemplo de tres etapas ($B=10$, $\sum t_i=22$) la fuerza bruta revisa los $2^3=8$ subconjuntos de la Tabla~\ref{tab:fb} y llega al mismo óptimo que la programación dinámica: acelerar las etapas 1 y 3.

\begin{table}[h]
\centering\small
\begin{tabular}{ccccc}
\toprule
Subconjunto & Peso = costo & Valor = ahorro & Duración $22-\text{ahorro}$ & ¿Cabe ($\le 10$)? \\
\midrule
$\emptyset$ & 0 & 0 & 22 & sí \\
$\{1\}$ & 4 & 3 & 19 & sí \\
$\{2\}$ & 3 & 2 & 20 & sí \\
$\{3\}$ & 5 & 4 & 18 & sí \\
$\{1,2\}$ & 7 & 5 & 17 & sí \\
\sh $\{1,3\}$ & \sh 9 & \sh \textbf{7} & \sh \textbf{15} & \sh sí (óptimo) \\
$\{2,3\}$ & 8 & 6 & 16 & sí \\
$\{1,2,3\}$ & 12 & 9 & 13 & no \\
\bottomrule
\end{tabular}
\caption{Los 8 subconjuntos que revisa la fuerza bruta. Las columnas sirven para la mochila (peso, valor) y para la aceleración (costo, ahorro, duración).}
\label{tab:fb}
\end{table}

\subsection{Programación dinámica: definición y recurrencia}

\textbf{Mochila.} $\DP{i}{b}$ es el máximo valor que se logra usando los primeros $i$ objetos con capacidad $b$.
\begin{equation}
  \DP{i}{b}=
  \begin{cases}
    \max\bigl(\DP{i-1}{b},\ \DP{i-1}{b-w_i}+v_i\bigr) & \text{si } w_i\le b,\\
    \DP{i-1}{b} & \text{si } w_i>b,
  \end{cases}
  \qquad \DP{0}{b}=0.
\end{equation}

\begin{algorithm}[h]
\caption{\textsc{Mochila-DP}$(w,v,W,n)$}
\begin{algorithmic}[1]
\For{$b\gets 0$ \textbf{hasta} $W$} \State $\DP{0}{b}\gets 0$ \EndFor
\For{$i\gets 1$ \textbf{hasta} $n$}
  \For{$b\gets 0$ \textbf{hasta} $W$}
    \State $\DP{i}{b}\gets\DP{i-1}{b}$
    \If{$w_i\le b$ \textbf{y} $\DP{i-1}{b-w_i}+v_i>\DP{i-1}{b}$}
      \State $\DP{i}{b}\gets\DP{i-1}{b-w_i}+v_i$
    \EndIf
  \EndFor
\EndFor
\State $b\gets W$;\quad $S^*\gets\emptyset$
\For{$i\gets n$ \textbf{hasta} $1$ \textbf{descendiendo}}
  \If{$\DP{i}{b}\ne\DP{i-1}{b}$} \State $S^*\gets S^*\cup\{i\}$;\quad $b\gets b-w_i$ \EndIf
\EndFor
\State \Return $(\DP{n}{W},S^*)$
\end{algorithmic}
\end{algorithm}

La reconstrucción no necesita guardar las decisiones en otra tabla. Como la comparación del \textbf{si} del llenado es estricta ($>$), un empate se resuelve como \textit{no tomar}; por eso $\DP{i}{b}\ne\DP{i-1}{b}$ ocurre exactamente cuando el objeto $i$ se tomó. Así está implementado en el código.

\textbf{Aceleración (equivalente).} $\DP{i}{b}$ es la máxima reducción alcanzable usando las primeras $i$ etapas y gastando como máximo $b$.
\begin{equation}
  \DP{i}{b}=
  \begin{cases}
    \max\bigl(\DP{i-1}{b},\ \DP{i-1}{b-c_i}+d_i\bigr) & \text{si } c_i\le b,\\
    \DP{i-1}{b} & \text{si } c_i>b,
  \end{cases}
  \qquad \DP{0}{b}=0.
\end{equation}
El algoritmo \textsc{Aceleración-DP} es el \textsc{Mochila-DP} con $(w,v,W)$ reemplazados por $(c,d,B)$, más un último paso que convierte la reducción máxima en la duración mínima de la cadena. Es el que está implementado en \texttt{src/aceleracion.py}.

\begin{algorithm}[h]
\caption{\textsc{Aceleración-DP}$(t,c,d,B,n)$}
\begin{algorithmic}[1]
\For{$b\gets 0$ \textbf{hasta} $B$} \State $\DP{0}{b}\gets 0$ \EndFor
\For{$i\gets 1$ \textbf{hasta} $n$}
  \For{$b\gets 0$ \textbf{hasta} $B$}
    \State $\DP{i}{b}\gets\DP{i-1}{b}$
    \If{$c_i\le b$ \textbf{y} $\DP{i-1}{b-c_i}+d_i>\DP{i-1}{b}$}
      \State $\DP{i}{b}\gets\DP{i-1}{b-c_i}+d_i$
    \EndIf
  \EndFor
\EndFor
\State $b\gets B$;\quad $A^*\gets\emptyset$
\For{$i\gets n$ \textbf{hasta} $1$ \textbf{descendiendo}}
  \If{$\DP{i}{b}\ne\DP{i-1}{b}$} \State $A^*\gets A^*\cup\{i\}$;\quad $b\gets b-c_i$ \EndIf
\EndFor
\State \Return $\bigl(\sum_i t_i-\DP{n}{B},\;A^*\bigr)$
\end{algorithmic}
\end{algorithm}

\textbf{Cómo leer la recurrencia.} La etapa $i$ solo tiene dos opciones:
\begin{itemize}
  \item \textit{No acelerarla:} el mejor valor es el de las $i-1$ etapas anteriores con el mismo $b$.
  \item \textit{Acelerarla:} se gana $d_i$ y a las etapas anteriores les queda un presupuesto de $b-c_i$.
\end{itemize}
Si la etapa no cabe en el presupuesto ($c_i>b$), solo queda la primera opción. La tabla toma el mejor de los dos casos, y $\DP{n}{B}$ es la reducción óptima.

\textbf{Principio de optimalidad.} La recurrencia funciona gracias al principio de optimalidad de Bellman (1957): una solución óptima de un problema está formada por soluciones óptimas de sus subproblemas. Aquí significa que, si una elección óptima con las primeras $i$ etapas y presupuesto $b$ acelera la etapa $i$, lo que eligió entre las $i-1$ etapas anteriores es a su vez óptimo para el presupuesto $b-c_i$; y si no la acelera, es óptimo para el presupuesto $b$. Si esa parte no fuera óptima, se podría cambiar por una mejor y se obtendría algo mejor que el óptimo, lo cual es imposible. Por eso basta con guardar en la fila $i-1$ el mejor valor de cada subproblema y combinar esos valores, sin revisar todas las combinaciones. En el ejemplo, $\DP{3}{10}=7$ se obtiene de $\DP{2}{5}=3$, que es el mejor ahorro posible con las etapas 1 y 2 y presupuesto 5.

\subsection{Ejemplo: las dos tablas}

Se usan los costos y ahorros del ejemplo: etapa 1 $=(c,d)=(4,3)$, etapa 2 $=(3,2)$, etapa 3 $=(5,4)$ y $B=10$. Para la mochila son los pesos y valores de los objetos y $W=10$.

\textbf{Tabla de los óptimos.} Cada casilla es el mayor valor (o ahorro) posible con las primeras $i$ filas y presupuesto $b$. La casilla de abajo a la derecha, $\DP{3}{10}=7$, es el máximo que se puede lograr con 10. Las casillas sombreadas son las que recorre la reconstrucción.

\begin{table}[h]
\centering\small
\setlength{\tabcolsep}{4.5pt}
\begin{tabular}{lrrrrrrrrrrr}
\toprule
$i\backslash b$ & 0 & 1 & 2 & 3 & 4 & 5 & 6 & 7 & 8 & 9 & 10 \\
\midrule
0 & 0 & \cellcolor{celda}\textbf{0} & 0 & 0 & 0 & 0 & 0 & 0 & 0 & 0 & 0 \\
1 & 0 & 0 & 0 & 0 & 3 & \cellcolor{celda}\textbf{3} & 3 & 3 & 3 & 3 & 3 \\
2 & 0 & 0 & 0 & 2 & 3 & \cellcolor{celda}\textbf{3} & 3 & 5 & 5 & 5 & 5 \\
3 & 0 & 0 & 0 & 2 & 3 & 4 & 4 & 5 & 6 & 7 & \cellcolor{celda}\textbf{7} \\
\midrule
$\sum t_i - \mathrm{DP}[3][b]$ & 22 & 22 & 22 & 20 & 19 & 18 & 18 & 17 & 16 & 15 & 15 \\
\bottomrule
\end{tabular}
\caption{Tabla $\DP{i}{b}$ con los valores óptimos: valor de la mochila o tiempo ahorrado en la aceleración. La última fila da la duración mínima de la cadena.}
\label{tab:moch-opt}
\end{table}

\textbf{Tabla de cómo se llegó.} Cada casilla muestra los dos candidatos \textit{no tomar}\,$|$\,\textit{tomar}, que son $\DP{i-1}{b}\,|\,\DP{i-1}{b-w_i}+v_i$. En negrita está el mayor, que es el valor de la tabla de óptimos. La flecha $\uparrow$ indica que el valor viene de la casilla de arriba (no se toma el objeto $i$) y $\nwarrow$ que viene de $(i-1,\,b-w_i)$ (se toma). El símbolo $-|-$ significa que no cabe ($w_i>b$). Un empate se resuelve como \textit{no tomar}.

\begin{table}[h]
\centering\footnotesize
\setlength{\tabcolsep}{1.7pt}
\begin{tabular}{lccccccccccc}
\toprule
$i\backslash b$ & 0 & 1 & 2 & 3 & 4 & 5 & 6 & 7 & 8 & 9 & 10 \\
\midrule
0 & 0 & 0 & 0 & 0 & 0 & 0 & 0 & 0 & 0 & 0 & 0 \\
1 & $-|-$\,$\uparrow$ & $-|-$\,$\uparrow$ & $-|-$\,$\uparrow$ & $-|-$\,$\uparrow$ & 0|\textbf{3}\,$\nwarrow$ & \cellcolor{celda}0|\textbf{3}\,$\nwarrow$ & 0|\textbf{3}\,$\nwarrow$ & 0|\textbf{3}\,$\nwarrow$ & 0|\textbf{3}\,$\nwarrow$ & 0|\textbf{3}\,$\nwarrow$ & 0|\textbf{3}\,$\nwarrow$ \\
2 & $-|-$\,$\uparrow$ & $-|-$\,$\uparrow$ & $-|-$\,$\uparrow$ & 0|\textbf{2}\,$\nwarrow$ & \textbf{3}|2\,$\uparrow$ & \cellcolor{celda}\textbf{3}|2\,$\uparrow$ & \textbf{3}|2\,$\uparrow$ & 3|\textbf{5}\,$\nwarrow$ & 3|\textbf{5}\,$\nwarrow$ & 3|\textbf{5}\,$\nwarrow$ & 3|\textbf{5}\,$\nwarrow$ \\
3 & $-|-$\,$\uparrow$ & $-|-$\,$\uparrow$ & $-|-$\,$\uparrow$ & $-|-$\,$\uparrow$ & $-|-$\,$\uparrow$ & 3|\textbf{4}\,$\nwarrow$ & 3|\textbf{4}\,$\nwarrow$ & \textbf{5}|4\,$\uparrow$ & 5|\textbf{6}\,$\nwarrow$ & 5|\textbf{7}\,$\nwarrow$ & \cellcolor{celda}5|\textbf{7}\,$\nwarrow$ \\
\bottomrule
\end{tabular}
\caption{Candidatos y decisión de cada casilla de la Tabla~\ref{tab:moch-opt}. En la mochila, $\uparrow$ es no tomar el objeto y $\nwarrow$ es tomarlo; en la aceleración, no acelerar o acelerar la etapa.}
\label{tab:moch-como}
\end{table}

\textbf{Cómo se lee.} Se parte de $(3,10)$: los candidatos son $5\,|\,7$, gana \textit{tomar}, así que el objeto 3 entra y la capacidad baja a $10-5=5$. En $(2,5)$ los candidatos son $3\,|\,2$, gana \textit{no tomar}. En $(1,5)$ son $0\,|\,3$, gana \textit{tomar}, y sobra $1$. Solución: $S^*=\{1,3\}$, peso $9\le 10$, valor $7$.

\textbf{Aceleración.} Las mismas dos tablas sirven, porque los números son los mismos: se leen con los nombres de la cadena. La última fila de la Tabla~\ref{tab:moch-opt} convierte el ahorro $\DP{3}{b}$ en la duración mínima, que con $B=10$ baja de $22$ a $15$. En la Tabla~\ref{tab:moch-como}, $\uparrow$ significa \textit{no acelerar la etapa $i$} y $\nwarrow$ significa \textit{acelerarla}.

\subsection{Reconstrucción: ¿qué etapas se aceleran?}

Se recorre la tabla hacia atrás desde $(n,B)$. Si $\DP{i}{b}\ne\DP{i-1}{b}$, la etapa $i$ se acelera y $b$ baja en $c_i$; si no, se conserva $b$. En el ejemplo, el recorrido (sección~\ref{sec:evid}) lleva a acelerar las etapas 1 y 3. Es el camino sombreado de las Tablas~\ref{tab:moch-opt} y~\ref{tab:moch-como}.

\subsection{Función de eficiencia}

\textbf{Fuerza bruta.} El ciclo interno hace $n$ iteraciones de costo constante $a$, y fuera de él hay un costo constante $b'$ por subconjunto (inicializar y comparar). Como el ciclo externo repite eso $2^n$ veces:
\begin{equation}
  T_{FB}(n)=\sum_{mascara=0}^{2^n-1}\bigl(a\,n+b'\bigr)=2^n\,(a\,n+b')=O(n\cdot 2^n).
\end{equation}
En la aceleración se suma además $O(n)$ para calcular $\sum t_i$, lo que no cambia la cota. La memoria extra es $O(1)$. Si en lugar de volver a sumar se decide objeto por objeto («no tomarlo» o «tomarlo») arrastrando el peso y el valor acumulados, la recurrencia es
\[
  T(n)=2\,T(n-1)+c,\qquad T(0)=d,
\]
que al desarrollarse da $T(n)=(d+c)\,2^n-c=O(2^n)$. Se gana el factor $n$, pero sigue siendo exponencial.

\textbf{Programación dinámica.} La inicialización escribe $(n+1)(B+1)$ casillas. Los dos ciclos evalúan $n(B+1)$ casillas con un número constante de operaciones cada una. La reconstrucción da $n$ pasos constantes. Con constantes $a$, $e$ y $r$:
\begin{equation}
  T_{DP}(n,B)=a\,(n+1)(B+1)+e\,n\,(B+1)+r\,n=O(n\cdot B).
\end{equation}
Esta última igualdad vale para $n\ge 1$ y $B\ge 1$. Para $n=0$ o $B=0$, el costo se expresa como $O\bigl((n+1)(B+1)\bigr)$. La memoria también es $O(n\cdot B)$, porque la reconstrucción necesita la tabla completa. En la mochila es $O(n\cdot W)$ por la misma cuenta.

La memoria se puede reducir: como cada fila solo depende de la anterior, basta guardar dos filas de valores, $O(B)$, si además se guarda en una tabla de un bit por casilla si el objeto se tomó, $O(n\cdot B)$ bits, para poder reconstruir. El orden sigue siendo $O(n\cdot B)$, pero con bits en lugar de enteros. Esta variante se comparará experimentalmente en una entrega posterior.

\begin{table}[h]
\centering\small
\begin{tabular}{lccc}
\toprule
Algoritmo & Tiempo & Memoria & Conviene cuando \\
\midrule
Fuerza bruta (iterativa) & $O(n\cdot 2^n)$ & $O(1)$ & $n$ muy pequeño \\
Fuerza bruta (recursiva) & $O(2^n)$ & $O(n)$ & $n$ muy pequeño \\
Programación dinámica & $O(n\cdot B)$ & $O(n\cdot B)$ & $B$ moderado \\
\bottomrule
\end{tabular}
\caption{Resumen de eficiencias. Para la mochila, $B$ se reemplaza por $W$.}
\end{table}

En el ejemplo la fuerza bruta evalúa $8$ subconjuntos y la programación dinámica $33$ casillas, así que con tamaños tan pequeños no hay diferencia. Con $n=30$ y $B=1\,000$ la fuerza bruta recursiva hace del orden de $2^{30}\approx 1{,}07\times 10^{9}$ pasos y la programación dinámica $30\cdot 1\,001\approx 3\times 10^{4}$.

\subsection{Por qué es pseudo-polinomial}

El tiempo crece con el valor de $B$, pero $B$ ocupa solo $k=\lceil\log_2(B+1)\rceil$ bits de la entrada. Por ejemplo, $B=1\,000\,000$ se escribe con 20 bits, pero la tabla tendría un millón de columnas. Reemplazando $B$ por su expresión en bits, el costo es exponencial en $k$:
\begin{equation}
  O(n\cdot B)=O(n\cdot 2^k).
\end{equation}
\textbf{Por qué no contradice la NP-dureza.} Un algoritmo polinomial tendría que serlo en la longitud $L$ de la entrada, y este no lo es. El problema es NP-difícil en sentido débil: es tratable cuando $B$ es pequeño, y se vuelve inmanejable cuando $B$ es enorme.

Esto también limita la ventaja sobre la fuerza bruta. Con $n=20$ y $B=10^{6}$ la tabla tiene unas $2\times 10^{7}$ casillas: tantas como los pasos de la fuerza bruta iterativa ($20\cdot 2^{20}\approx 2\times 10^{7}$) y más que los de la recursiva ($2^{20}\approx 10^{6}$). Frente a la recursiva, $O(2^n)$, la programación dinámica gana mientras $n\cdot B$ sea menor que $2^n$; frente a la iterativa, $O(n\cdot 2^n)$, mientras $B$ sea menor que $2^n$.

% ======================= 6 EVIDENCIA =======================
\section{Evidencia de ejecución}
\label{sec:evid}

\textbf{Qué se probó.} La solución pasó las 19 pruebas automatizadas, incluidos 300 casos aleatorios comparados contra fuerza bruta (probar todas las combinaciones). Los casos de evaluación del curso no se comparten, así que el grupo creó sus propios casos, y su respuesta esperada se calculó por fuerza bruta, sin usar la programación dinámica. El monitor compartió un caso de referencia de 20 etapas en un formato que incluye también datos de auditoría y de variantes de entregas posteriores. Como esta entrega cubre solo la aceleración, se mantuvo el formato propio y los casos del grupo.

\textbf{Comprobación por fuerza bruta de cada caso.} Además de las pruebas, cada caso de hasta 20 etapas se comprueba contra la fuerza bruta: el programa revisa una por una las $2^n$ combinaciones de etapas, registra de cada una el costo, si cabe en el presupuesto y el ahorro, y verifica que la programación dinámica haya elegido una de las soluciones óptimas. Los 9 casos que se pueden comprobar coinciden; el de 200 etapas no se comprueba, porque tendría $2^{200}$ combinaciones. El registro de cada caso queda en \texttt{resultados/comprobacion/<id>.log}.

\textbf{Resultados por caso.} La Tabla~\ref{tab:resultados} resume el resultado de cada caso propio: qué prueba, su tamaño, las etapas aceleradas, el costo usado, la duración antes y después, y el tiempo. El tiempo es el promedio de 5 ejecuciones y cambia ligeramente en cada corrida.

\begin{table}[h]
\centering\scriptsize
\setlength{\tabcolsep}{2.6pt}
\begin{tabular}{@{}>{\ttfamily}p{3.2cm}p{3.5cm}rrcrcr@{}}
\toprule
\normalfont Caso & Qué prueba & $n$ & $B$ & Aceleradas & Costo & Duración & Tiempo (ms) \\
\midrule
ejemplo\_tres\_etapas & Caso básico & 3 & 10 & 1, 3 & 9 & $22\to 15$ & 0,007 \\
ejemplo\_almuerzo & Presupuesto usado completo & 3 & 5 & 2, 3 & 5 & $28\to 6$ & 0,004 \\
voraz\_por\_razon\_falla & Elegir por $d/c$ no es óptimo & 3 & 50 & 2, 3 & 50 & $280\to 60$ & 0,011 \\
borde\_presupuesto\_cero & $B=0$ & 2 & 0 & ninguna & 0 & $9\to 9$ & 0,002 \\
borde\_costo\_cero & Etapa con costo 0 & 2 & 0 & 1 & 0 & $9\to 7$ & 0,003 \\
borde\_nada\_cabe & Ninguna etapa cabe & 2 & 2 & ninguna & 0 & $9\to 9$ & 0,003 \\
borde\_todo\_cabe & Todas las etapas caben & 3 & 100 & 1, 2, 3 & 6 & $16\to 5$ & 0,025 \\
borde\_sin\_reduccion & Etapa con $d=0$ & 2 & 5 & 2 & 2 & $9\to 5$ & 0,003 \\
aleatorio\_n15\_s7 & Caso generado, $n=15$ & 15 & 33 & 11 etapas & 29 & $156\to 69$ & 0,041 \\
aleatorio\_n200\_s1 & Caso grande, $n=200$ & 200 & 824 & 113 etapas & 824 & $2130\to 1188$ & 13,204 \\
\bottomrule
\end{tabular}
\caption{Resultados de la solución base por caso de prueba. \textit{Aceleradas}: etapas que se aceleran; \textit{Costo}: costo usado; \textit{Duración}: original $\to$ final.}
\label{tab:resultados}
\end{table}

Los casos de borde cubren situaciones límite: presupuesto cero, etapas gratis, ninguna etapa que quepa y todas las etapas aceleradas a la vez. El caso \texttt{voraz\_por\_razon\_falla} muestra que elegir siempre la etapa con mejor razón $d/c$ no siempre es óptimo, y por eso se usa programación dinámica.

\textbf{Traza de cada caso.} Cada caso tiene su archivo \texttt{resultados/<id>.json} en el repositorio, con el mismo nombre del caso. Además del resultado y el tiempo, guarda la traza: cada paso de la reconstrucción ($i$, $b$, $\DP{i}{b}$, $\DP{i-1}{b}$, la decisión y el presupuesto restante) y, cuando tiene menos de 500 casillas, la tabla DP completa. La subsección~\ref{sub:traza} muestra la traza del caso \texttt{ejemplo\_tres\_etapas}.

La guía de la entrega también pide reportar la cobertura de auditoría. Esa es la segunda parte del Subproyecto~3, la cobertura mínima de auditoría, y no aplica en este informe, que cubre solo la aceleración de la cadena (sección~\ref{sec:alcance}).

\subsection{Traza del caso \texttt{ejemplo\_tres\_etapas}}
\label{sub:traza}

Las etapas $(t,c,d)$ son: $1=(8,4,3)$, $2=(5,3,2)$, $3=(9,5,4)$, con $B=10$. La reducción total es 7, con costo 9, y la duración baja de 22 a 15. Así se reconstruye la respuesta recorriendo la tabla de la sección~\ref{sec:sol} hacia atrás:

\begin{table}[h]
\centering
\begin{tabular}{cccccc}
\toprule
Etapa $i$ & $b$ & $\DP{i}{b}$ & $\DP{i-1}{b}$ & Decisión & Presupuesto restante \\
\midrule
3 & 10 & 7 & 5 & Acelerar & 5 \\
2 & 5 & 3 & 3 & No acelerar & 5 \\
1 & 5 & 3 & 0 & Acelerar & 1 \\
\bottomrule
\end{tabular}
\caption{Traza de la reconstrucción.}
\label{tab:traza}
\end{table}

Cómo se lee: en la etapa 3, $\DP{3}{10}=7$ es distinto de $\DP{2}{10}=5$, así que se acelera y el presupuesto baja de 10 a 5. En la etapa 2 los dos valores coinciden (3), así que no se acelera. En la etapa 1, 3 es distinto de 0, así que se acelera.

\section*{Referencias}
\addcontentsline{toc}{section}{Referencias}
\begin{enumerate}[label={[\arabic*]},leftmargin=2em,itemsep=4pt]
  \item Bellman, R. (1957). \textit{Dynamic Programming}. Princeton University Press.
  \item Cook, S. A. (1971). The complexity of theorem-proving procedures. En \textit{Proceedings of the Third Annual ACM Symposium on Theory of Computing} (pp. 151--158).
  \item Karp, R. M. (1972). Reducibility among combinatorial problems. En \textit{Complexity of Computer Computations} (pp. 85--103). Plenum Press.
  \item \textit{Proyecto Semestral, Análisis y Diseño de Algoritmos, 2026-B}. Universidad de Caldas.
  \item \textit{Guía de requisitos, primera entrega, ADAG 2026-B}.
\end{enumerate}

\end{document}
